\section*{Step 250: Branch C Higher-\(k\) Growth Diagnostic}

\paragraph{Kernel derivative.}
Burnol's 2002 formula gives
\[
K_a^\Gamma(z,w)=
\frac{E_\lambda(z)\overline{E_\lambda(w)}
-E_\lambda(1-z)\overline{E_\lambda(1-w)}}{z+w-1}.
\]
Since \(z+w-1\) is independent of \(\overline w\), the higher
anti-holomorphic derivatives have the form
\[
\partial_{\overline w}^{k}K_a^\Gamma(z,w)
=
\frac{
E_\lambda(z)\overline{E_\lambda^{(k)}(w)}
 + (-1)^{k+1}E_\lambda(1-z)\overline{E_\lambda^{(k)}(1-w)}
}{z+w-1}.
\]
In the critical-line numerical realization this is implemented as
\[
L_{\rho,k}(G)
=(-i\partial_\gamma)^k
\bigl(P_\infty M_\zeta G\bigr)(1/2+i\gamma)\big|_{\gamma=\Im\rho}.
\]

\paragraph{Computed values.}
For the triples \((\rho_1,G_\star)\), \((\rho_2,G_\star)\), and
\((\rho_1,G')\), the new values are:
\[
\begin{array}{c|cc}
\text{triple} & k=3 & k=4\\
\hline
\rho_1,G_\star & 1.1164964359 & 2.2021731761\\
\rho_2,G_\star & 1.1457543044 & 2.5868233085\\
\rho_1,G' & 1.3778472859 & 2.5398826197
\end{array}
\]
with numerical error bounds below \(1.2\cdot10^{-5}\).

\paragraph{Growth-law fit.}
Linear, quadratic, exponential, factorial, affine-factorial, and
Pochhammer models were fitted over \(k=0,\ldots,4\).  The exponential model
had the lowest raw RMSE for all three triples:
\[
\mathrm{RMSE}=2.74\cdot10^{-3},\quad
3.38\cdot10^{-2},\quad
3.50\cdot10^{-3}.
\]

\paragraph{Verdict.}
The finite sampled diagnostic supports
\[
\mathrm{V\_branch\_C\_k\_growth\_law\_exponential}.
\]
This is not an asymptotic theorem and does not claim RH closure.
